\MathBookSourcePage{313}
\subsection{习题}
\label{sec:ch10-exercises}
\setcounter{exerciseitem}{0}

\begin{MBExercise}
\label{exr:ch10-01}
证明式~\eqref{eq:ch10-lobatto-polynomial} 定义的多项式 $L_k$ 有互异零点 $0=\xi_0<\xi_1<\cdots<\xi_{k-1}=1$. 提示: 设 $\xi_1<\cdots<\xi_r$ 是 $(0,1)$ 中奇重数零点. 证明 $L_k(x)\prod_{i=1}^r(x-\xi_i)$ 在 $[0,1]$ 上不变号; 再利用式~\eqref{eq:ch10-lobatto-remainder-vanishes} 证明 $r<k-2$ 会导致矛盾.
\end{MBExercise}

\begin{MBExercise}
\label{exr:ch10-02}
证明推论~\ref{cor:ch10-lobatto-error}. 提示: 加上再减去 $Q^{2k-2}f$, 并使用第~\MBUnresolvedReference{sec:ch04-approximation}{4.4} 节的方法.
\end{MBExercise}

\begin{MBExercise}
\label{exr:ch10-03}
证明当 $h$ 充分小时, 曲边三角形上的 Lagrange 元良定义. 提示: 用伸缩映射到一族尺度为一且有一条曲边的参考单元. 注意该曲边到一条固定直边的距离仅为 $O(h)$, 再作扰动论证.
\end{MBExercise}

\begin{MBExercise}
\label{exr:ch10-04}
证明当 $h$ 充分小时, 定理~\ref{thm:ch10-interpolated-boundary-error} 可改进为
\[
\MathBookFormulaID{formula-ch10-ex04-improved-boundary-error}
\|u-u_h\|_{H^1(\Omega)}\leq C_\rho h^{k-1}\|u\|_{H^k(\Omega)}.
\]
提示: 用多项式 $P(s)$ 逼近 $\partial u/\partial\nu$, 并分别估计
\[
\MathBookFormulaID{formula-ch10-ex04-boundary-split}
\int_{\partial\Omega}\left(\frac{\partial u}{\partial\nu}-P\right)w\,ds,
\qquad
\int_{\partial\Omega}Pw\,ds.
\]
使用当 $\|w\|_{W_\infty^1(\Omega)}=1$ 时 $w=O(h^2)$ 于 $\partial\Omega$ 上这一事实.
\end{MBExercise}

\begin{MBExercise}
\label{exr:ch10-05}
证明引理~\ref{lem:ch10-symmetric-strang}. 提示: 取 $v$ 为 $u$ 到 $V_h$ 上的正交投影, 并仿照式~\eqref{eq:ch10-first-strang-proof-chain}.
\end{MBExercise}

\begin{MBExercise}
\label{exr:ch10-06}
定理~\ref{thm:ch10-isoparametric-error} 中关于 $f$ 的范数能否改进? 例如, 当 $k\leq3+n/2$ 时, 是否可以得到
\[
\MathBookFormulaID{formula-ch10-ex06-isoparametric-improved-load}
\|u-\widehat u_h\|_{H^1(\Omega)}
\leq C_\rho h^{k-1}\bigl(\|u\|_{H^k(\Omega)}+\|f\|_{H^{k-2}(\Omega)}\bigr)?
\]
提示: 考虑 $\|f-f\circ\Phi^h\|_{H^{-1}(\Omega)}$ 的界.
\end{MBExercise}

\begin{MBExercise}
\label{exr:ch10-07}
证明迹估计~\eqref{eq:ch10-scaled-trace-edge}. 提示: 在参考单纯形上使用迹定理, 再作齐次性/尺度变换论证.
\end{MBExercise}

\begin{MBExercise}
\label{exr:ch10-08}
证明离散估计~\eqref{eq:ch10-jump-midpoint-estimate}.
\end{MBExercise}

\begin{MBExercise}
\label{exr:ch10-09}
证明图~\ref{fig:ch10-morley-element} 所示的 Morley 有限元是有限元. 这里 $K$ 是三角形, $\mathcal P$ 是二次多项式空间, $\mathcal N$ 由顶点处的函数值和边中点处的法向导数值组成. 提示: 在三角形顶点为零的二次多项式在边中点达到极值. 用此证明对 $v\in\mathcal P$, $\mathcal Nv=\{0\}$ 蕴含 $\nabla v\equiv0$.
\end{MBExercise}

\MathBookSourcePage{314}
\begin{MBExercise}
\label{exr:ch10-10}
考虑第~\MBUnresolvedReference{sec:ch05-biharmonic}{5.9} 节的双调和问题, 定义非协调双线性形式
\[
\MathBookFormulaID{formula-ch10-ex10-morley-bilinear-form}
a_h(u,v)=\sum_{T\in\mathcal T^h}\int_T
\bigl[\Delta u\Delta v-(1-\nu)(2u_{xx}v_{yy}+2u_{yy}v_{xx}-4u_{xy}v_{xy})\bigr],dxdy,
\]
其中 $\mathcal T^h$ 是 $\Omega$ 的三角剖分, $0<\nu<1$. 令 $V_h$ 是与 $\mathcal T^h$ 相联系的 Morley 有限元空间. 证明 $a_h(\cdot,\cdot)$ 在 $V_h\cup V$ 上非退化. 提示: 使用式~\MBUnresolvedReference{eq:ch05-source-5-9-2}{5.9.2} 以及式~\eqref{eq:ch10-broken-energy-norm} 后面的讨论.
\end{MBExercise}

\begin{MBExercise}
\label{exr:ch10-11}
令 $G$ 是具有公共边 $e$ 的两个非退化三角形之并, $\operatorname{diam}G=1$. 令 $V$ 由满足下列条件的 $z$ 组成: $z|_{T_i}$ 是二次多项式, $z$ 在 $e$ 的端点连续, $\nabla z$ 在 $e$ 的中点连续. 证明存在 $C<\infty$, 使得对所有 $z\in V$,
\[
\MathBookFormulaID{formula-ch10-ex11-piecewise-quadratic-approximation}
\inf_{P\in\mathcal P_1}\|z-P\|_{L^2(G)}
\leq C\sum_{i=1}^2|z|_{H^2(T_i)}.
\]
提示: 对在 $T$ 一条边的顶点为零且在该边中点的法向导数为零的所有 $z\in\mathcal P_2$, 证明 $\|z\|_{L^2(T)}\leq C|z|_{H^2(T)}$.
\end{MBExercise}

\begin{MBExercise}
\label{exr:ch10-12}
令 $D^kf$, $1\leq k\leq3$, 表示由分部积分公式
\[
\MathBookFormulaID{formula-ch10-ex12-biharmonic-integration-by-parts}
\int_T[\Delta z\Delta\zeta-(1-\nu)(2z_{xx}\zeta_{yy}+2z_{yy}\zeta_{xx}-4z_{xy}\zeta_{xy})],dxdy
=\int_{\partial T}(D^1zD^2\zeta+zD^3\zeta)\,ds+\int_Tz\Delta^2\zeta\,dxdy
\]
产生的 $f$ 的 $k$ 阶导数. 令 $G$ 和 $V$ 如习题~\ref{exr:ch10-11}. 证明存在 $C<\infty$, 使得
\[
\MathBookFormulaID{formula-ch10-ex12-edge-boundary-term}
\left|\int_eD^1(z|_{T_1}-z|_{T_2})D^2\zeta\,ds
+\int_e(z|_{T_1}-z|_{T_2})D^3\zeta\,ds\right|
\leq C|z|_G\bigl(|\zeta|_{H^3(G)}+\|\Delta^2\zeta\|_{L^2(G)}\bigr),
\]
其中 $z\in V$, $\zeta\in H^3(G)$, $\Delta^2\zeta\in L^2(G)$, 且
\[
\MathBookFormulaID{formula-ch10-ex12-z-energy}
|z|_G^2=\sum_{i=1}^2\int_{T_i}\bigl[(\Delta z)^2-(1-\nu)(4z_{xx}z_{yy}-4z_{xy}^2)\bigr],dxdy.
\]
提示: 用分部积分公式证明边界项由 $|z|_G|\zeta|_G+\|z\|_{L^2(G)}\|\Delta^2\zeta\|_{L^2(G)}$ 控制, 再应用 Bramble--Hilbert 引理~\MBUnresolvedReference{lem:ch04-source-4-3-8}{4.3.8} 和习题~\ref{exr:ch10-11}. 利用向 $\zeta$ 加二次多项式、向 $z$ 加任意光滑函数都不改变该边界项的事实.
\end{MBExercise}

\MathBookSourcePage{315}
\begin{MBExercise}
\label{exr:ch10-13}
令 $a_h(\cdot,\cdot)$ 和 $V_h$ 如习题~\ref{exr:ch10-10}, $u_h\in V_h$ 满足
\[
\MathBookFormulaID{formula-ch10-ex13-morley-discrete-problem}
a_h(u_h,v)=\int_\Omega fv\,dx\qquad\forall v\in V_h.
\]
假设式~\MBUnresolvedReference{eq:ch05-source-5-9-4}{5.9.4} 的数据和解满足 $f\in L^2(\Omega)$, $u\in H^3(\Omega)$. 证明
\[
\MathBookFormulaID{formula-ch10-ex13-morley-error}
\|u-u_h\|_h\leq Ch\bigl(|u|_{H^3(\Omega)}+\|f\|_{L^2(\Omega)}\bigr),
\]
其中 $\|u\|_h:=\sqrt{a_h(u,u)}$. 提示: 使用习题~\ref{exr:ch10-12}.
\end{MBExercise}

\begin{MBExercise}
\label{exr:ch10-14}
考虑习题~\ref{exr:ch10-13} 所述基于 Morley 元的双调和问题非协调方法. 只假设式~\MBUnresolvedReference{eq:ch05-source-5-9-4}{5.9.4} 的解 $u\in H^3(\Omega)$, 能否在不假设 $f\in L^2(\Omega)$ 的条件下证明 $\|u-u_h\|_h\leq Ch|u|_{H^3(\Omega)}$?
\end{MBExercise}

\begin{MBExercise}
\label{exr:ch10-15}
像习题~\MBUnresolvedReference{exr:ch00-source-0-x-11}{0.x.11} 那样用数值求积逼近右端是最简单的变分罪行. 假设变分问题如第~\MBUnresolvedReference{sec:ch05-duality}{5.4} 节, 即 $a(u_h,v)=(f,v)$, $\forall v\in V_h$, 其中 $V_h$ 由次数不超过 $r$ 的分片多项式组成; 令 $\widetilde u_h\in V_h$ 满足
\[
\MathBookFormulaID{formula-ch10-ex15-quadrature-discrete-problem}
a(\widetilde u_h,v)=Q(fv)\qquad\forall v\in V_h,
\]
其中求积逼近满足
\[
\MathBookFormulaID{formula-ch10-ex15-quadrature-assumption}
\left|\int_\Omega v(x)w(x)\,dx-Q(vw)\right|
\leq Ch^k\sum_{T\in\mathcal T^h}\|v\|_{H^k(T)}\|w\|_{H^k(T)}.
\]
证明
\[
\MathBookFormulaID{formula-ch10-ex15-quadrature-error}
\|u_h-\widetilde u_h\|_{H^1(\Omega)}
\leq Ch^{k-r+1}\|f\|_{H^k(\Omega)}.
\]
提示: 见习题~\MBUnresolvedReference{exr:ch00-source-0-x-13}{0.x.13} 和~\MBUnresolvedReference{exr:ch00-source-0-x-14}{0.x.14}, 并使用逆估计.
\end{MBExercise}

\MathBookSourcePage{316}
\begin{MBExercise}
\label{exr:ch10-16}
除习题~\ref{exr:ch10-15} 的求积规则假设外, 再假设 $Q(w)=\sum_{T\in\mathcal T^h}Q_T(w)$, 且 $|Q_T(w)|\leq C\operatorname{meas}(T)\|w\|_{L^\infty(T)}$. 当 $k-r+1>n/2$ 时, 把习题~\ref{exr:ch10-15} 的估计改进为
\[
\MathBookFormulaID{formula-ch10-ex16-improved-quadrature-error}
\|u_h-\widetilde u_h\|_{H^1(\Omega)}
\leq Ch^{k-r+1}\|f\|_{H^{k-r+1}(\Omega)}.
\]
提示: 从 $f$ 中减去其插值函数.
\end{MBExercise}

\begin{MBExercise}
\label{exr:ch10-17}
在变系数情形(见第~\MBUnresolvedReference{sec:ch05-variable-coefficients}{5.7} 节)用数值求积逼近变分形式会导致引理~\ref{lem:ch10-second-strang} 所述的变分罪行. 令 $a(\cdot,\cdot)$ 如式~\MBUnresolvedReference{eq:ch05-source-5-6-3}{5.6.3}, 并定义
\[
\MathBookFormulaID{formula-ch10-ex17-quadrature-form}
a_h(u,v):=Q\left(\sum_{i,j=1}^na_{ij}(x)
\frac{\partial u}{\partial x_i}(x)\frac{\partial v}{\partial x_j}(x)\right),
\]
其中 $Q$ 满足习题~\ref{exr:ch10-15} 的条件. 假设 $V_h$ 由次数不超过 $r$ 的分片多项式组成, $a(\cdot,\cdot)$ 在 $V$ 上强制, 且 $k>2r-2$. 证明 $a_h(\cdot,\cdot)$ 在 $V_h$ 上强制. 提示: 用逆估计证明
\[
\MathBookFormulaID{formula-ch10-ex17-form-perturbation}
|a(v,v)-a_h(v,v)|\leq Ch^{k-2r+2}\|v\|_{H^1(\Omega)}^2
\qquad\forall v\in V_h.
\]
注意无需假设求积权为正.
\end{MBExercise}

\begin{MBExercise}
\label{exr:ch10-18}
在习题~\ref{exr:ch10-17} 的假设下, 定义 $\widetilde u_h\in V_h$ 满足 $a_h(\widetilde u_h,v)=(f,v)$, $\forall v\in V_h$. 证明
\[
\MathBookFormulaID{formula-ch10-ex18-variable-coefficient-error}
\|u-\widetilde u_h\|_{H^1(\Omega)}
\leq Ch^r\|u\|_{H^{k+1}(\Omega)}
\left(1+\max_{ij}\|a_{ij}\|_{W_\infty^k(\Omega)}\right).
\]
提示: 见习题~\MBUnresolvedReference{exr:ch00-source-0-x-13}{0.x.13} 和~\MBUnresolvedReference{exr:ch00-source-0-x-14}{0.x.14}, 应用引理~\ref{lem:ch10-second-strang} 并使用逆估计.
\end{MBExercise}

\begin{MBExercise}
\label{exr:ch10-19}
在习题~\ref{exr:ch10-16} 和~\ref{exr:ch10-17} 的假设下, 改进习题~\ref{exr:ch10-18} 中关于 $u$ 和 $a_{ij}$ 的范数. 提示: 从 $u$ 和 $a_{ij}$ 中减去相应插值函数.
\end{MBExercise}

\begin{MBExercise}
\label{exr:ch10-20}
考虑第~\MBUnresolvedReference{sec:ch05-duality}{5.4} 节所述 Poisson 方程的变分逼近, 在单纯形网格 $\mathcal T^h$ 上使用分片线性函数, 但像习题~\ref{exr:ch10-15} 那样对右端作求积逼近. 具体地, 令 $Q(w)=\sum_{T\in\mathcal T^h}Q_T(w)$, 其中
\[
\MathBookFormulaID{formula-ch10-ex20-trapezoidal-rule}
Q_T(w):=\operatorname{meas}(T)\sum_{i=1}^nw(z_i^T)
\qquad\forall T\in\mathcal T^h,
\]
$\{z_i^T:i=1,\ldots,n\}$ 是 $T$ 的顶点(这是梯形公式的 $n$ 维版本). 证明在 $\Omega=[0,1]\times[0,1]$ 上由 $45^\circ$ 直角三角形组成的规则网格(参见第~\MBUnresolvedReference{sec:ch00-difference-method}{0.5} 节), 相应的内部``差分模板''是标准五点差分模板. 分别描述 Neumann 和 Dirichlet 边界条件下的边界模板, 并计算三维规则网格上的差分模板.
\end{MBExercise}

\MathBookSourcePage{317}
\begin{MBExercise}
\label{exr:ch10-21}
证明习题~\ref{exr:ch10-20} 定义的 $n$ 维梯形公式满足习题~\ref{exr:ch10-16} 的条件, 并且当 $n\leq3$ 时,
\[
\MathBookFormulaID{formula-ch10-ex21-trapezoidal-product-error}
\left|\int_Tv(x)w(x)\,dx-Q_T(vw)\right|
\leq Ch^2\|v\|_{H^2(T)}\|w\|_{H^2(T)},
\]
其中 $h:=\operatorname{diam}T$. 提示: 证明 $Q_T$ 对线性函数精确, 再使用 Sobolev 不等式和第~\ref{chap:ch04-polynomial-approximation} 章的逼近结果.
\end{MBExercise}

\begin{MBExercise}
\label{exr:ch10-22}
可用分片线性函数和下述网格上的梯形公式产生 Poisson 方程差分方法~\MBUnresolvedReference{eq:ch00-source-0-5-3}{0.5.3} 的二维推广(另见习题~\MBUnresolvedReference{exr:ch00-source-0-x-11}{0.x.11}). 令 $0=x_0<x_1<\cdots<x_m=1$, $0=y_0<y_1<\cdots<y_n=1$ 为 $[0,1]$ 的分割, 以顶点 $\{(x_i,y_j):0\leq i\leq m,0\leq j\leq n\}$ 构造三角网格, 例如令所有边平行于坐标轴或沿 $(x_i,y_j)$ 到 $(x_{i+1},y_{j+1})$ 的对角线. 证明在适当的 $u,f$ 假设下,
\[
\MathBookFormulaID{formula-ch10-ex22-maximum-norm-difference-error}
\max_{i,j}|u(x_i,y_j)-U_{ij}|\leq Ch^2|\log h|.
\]
提示: 使用习题~\ref{exr:ch10-18}、\ref{exr:ch10-20} 和~\ref{exr:ch10-21}, 以及分片线性函数的 $\|v\|_{L^\infty(\Omega)}\leq C|\log h|\|v\|_{H^1(\Omega)}$. 另见第~\MBUnresolvedReference{sec:ch08-source-8-5}{8.5} 节和 Scott (1976).
\end{MBExercise}

\begin{MBExercise}
\label{exr:ch10-23}
考虑习题~\MBUnresolvedReference{exr:ch06-source-6-x-11}{6.x.11} 中的求积公式. 对哪些分片多项式次数, 它能为变系数问题给出最佳阶逼近? 陈述并证明适当的收敛定理. 提示: 使用习题~\ref{exr:ch10-18}.
\end{MBExercise}

\begin{MBExercise}
\label{exr:ch10-24}
陈述并证明一个同时考虑变系数和右端求积效应的收敛定理, 即 $\widetilde u_h\in V_h$ 由
\[
\MathBookFormulaID{formula-ch10-ex24-combined-quadrature-problem}
a_h(\widetilde u_h,v)=Q(fv)\qquad\forall v\in V_h
\]
定义.
\end{MBExercise}

\MathBookSourcePage{318}
\begin{MBExercise}
\label{exr:ch10-25}
令 $V_h$ 是与式~\eqref{eq:ch10-nonconforming-p1-space} 中三角剖分 $\mathcal T^h$ 相联系的非协调 $\mathcal P_1$ 有限元空间, $\widetilde V_h\subset H_0^1(\Omega)$ 是与 $\mathcal T^h$ 相联系的 $\mathcal P_2$ Lagrange 有限元空间. 定义 $E_h:V_h\to\widetilde V_h$:

(i) 在任一中点 $m$, $(E_hv)(m)=v(m)$;

(ii) 在任一顶点 $p$, $(E_hv)(p)$ 等于 $p$ 相邻中点处 $v$ 的平均值.

定义 $F_h:\widetilde V_h\to V_h$ 为 $(F_h\widetilde v)(m)=\widetilde v(m)$. 证明:

(a) 对所有 $v\in V_h$, $F_hE_hv=v$;

(b) 对所有 $v\in V_h$, $\|v-E_hv\|_{L^2(\Omega)}\leq Ch\|v\|_h$;

(c) 对所有 $\widetilde v\in\widetilde V_h$, $\|\widetilde v-F_h\widetilde v\|_{L^2(\Omega)}\leq Ch|\widetilde v|_{H^1(\Omega)}$.
\end{MBExercise}

\begin{MBExercise}
\label{exr:ch10-26}
使用习题~\ref{exr:ch10-25} 的结果证明非协调 $\mathcal P_1$ 有限元的 Poincaré 不等式 $\|v\|_{L^2(\Omega)}\leq C\|v\|_h$, $\forall v\in V_h$. 提示: 使用 Poincaré 不等式(参见命题~\MBUnresolvedReference{prop:ch05-source-5-3-5}{5.3.5})和逆估计(参见定理~\MBUnresolvedReference{thm:ch04-source-4-5-11}{4.5.11} 及备注~\MBUnresolvedReference{rem:ch04-source-4-5-20}{4.5.20}).
\end{MBExercise}

\begin{MBExercise}
\label{exr:ch10-27}
$\mathcal P_2$ Lagrange 有限元是非协调 $\mathcal P_1$ 有限元的协调亲属(Brenner 1996a), 意思是后者的形函数和节点变量也是前者的形函数和节点变量. 为旋转 $Q_1$ 元(参见习题~\MBUnresolvedReference{exr:ch03-source-3-x-15}{3.x.15})找到一个协调亲属, 并构造具有类似性质的 $E_h,F_h$. 推导旋转 $Q_1$ 元的 Poincaré 不等式.
\end{MBExercise}

\begin{MBExercise}
\label{exr:ch10-28}
为 Morley 有限元找到一个协调亲属(参见习题~\ref{exr:ch10-27}).
\end{MBExercise}

\begin{MBExercise}
\label{exr:ch10-29}
令 $\mathcal P^h$ 是有界多边形区域 $\Omega\subset\mathbb R^2$ 的分割族, $\mathcal T^h$ 是把 $\mathcal P^h$ 中每个三角形的中心与属于该三角形边界的 $\mathcal V^h$ 中各点连接所得的三角剖分族. 证明 $\mathcal P^h$ 非退化当且仅当 $\mathcal T^h$ 非退化.
\end{MBExercise}

\begin{MBExercise}
\label{exr:ch10-30}
在 $\Omega$ 为有界多边形区域的假设下证明式~\eqref{eq:ch10-element-green-formula}. 椭圆正则性说明 $u$ 在凹角以外属于 $H^2$, 因而只需考虑三角形 $T$ 的顶点 $p$ 是 $\Omega$ 的内角 $\omega>\pi$ 的凹角. 椭圆正则性理论(参见 Grisvard 1985, Dauge 1988, 或 Nazarov and Plamenevsky 1994)在 $p$ 附近给出
\[
\MathBookFormulaID{formula-ch10-ex30-corner-singularity}
u=u_R+\kappa r^{\pi/\omega}\sin((\pi/\omega)\theta),
\]
其中 $u_R\in H^2(\Omega)$, $\kappa\in\mathbb R$, $(r,\theta)$ 是 $p$ 处的极坐标, 从 $p$ 出发的两条 $\Omega$ 边分别由 $\theta=0$ 和 $\theta=\omega$ 给出. 令 $T_\delta=\{x\in T:|x-p|>\delta\}$, 在
\[
\MathBookFormulaID{formula-ch10-ex30-truncated-green}
\int_{T_\delta}\nabla u\cdot\nabla v\,dx
=\int_{T_\delta}fv\,dx+\int_{\partial T_\delta}\nabla u\cdot vn\,ds
\]
中令 $\delta\downarrow0$, 从而证明式~\eqref{eq:ch10-element-green-formula}.
\end{MBExercise}

\MathBookSourcePage{319}
\begin{MBExercise}
\label{exr:ch10-31}
证明式~\eqref{eq:ch10-ip-energy-norm} 和~\eqref{eq:ch10-ip-triple-norm} 定义的 $\|\cdot\|_h$ 与 $|||\cdot|||_h$ 是 $\mathring H^1(\Omega)+V_h$ 上的范数.
\end{MBExercise}

\begin{MBExercise}
\label{exr:ch10-32}
使用 $L^2(e)$ 上的 Cauchy--Schwarz 不等式和离散 Cauchy--Schwarz 不等式验证式~\eqref{eq:ch10-ip-boundedness}.
\end{MBExercise}

\begin{MBExercise}
\label{exr:ch10-33}
解释为什么第~\ref{sec:ch10-departure-framework} 节的结果不能用于第~\ref{sec:ch10-discontinuous} 节的内部罚方法.
\end{MBExercise}

\begin{MBExercise}
\label{exr:ch10-34}
检查式~\eqref{eq:ch10-ec-wp-difference} 中的常数与 $p$ 无关.
\end{MBExercise}

% The source has no Exercise 10.x.35.
\setcounter{exerciseitem}{35}
\begin{MBExercise}
\label{exr:ch10-36}
补充定理~\ref{thm:ch10-piecewise-poincare-infinity} 的证明细节.
\end{MBExercise}
